\documentclass[11pt,a4paper]{article}
\usepackage{fontspec,amsmath,amssymb,amsthm,microtype}
\setmainfont{lmroman10-regular.otf}[BoldFont=lmroman10-bold.otf,ItalicFont=lmroman10-italic.otf,BoldItalicFont=lmroman10-bolditalic.otf]
\usepackage{xcolor}
\usepackage[most]{tcolorbox}
\usepackage{mathrsfs}
\providecommand{\lt}{<}
\providecommand{\gt}{>}
\setmonofont{DejaVuSansMono.ttf}[BoldFont=DejaVuSansMono-Bold.ttf,ItalicFont=DejaVuSansMono-Oblique.ttf,BoldItalicFont=DejaVuSansMono-BoldOblique.ttf,Scale=0.86]
\newfontfamily\leanfallback{STIXTwoMath-Regular.otf}[Scale=0.86]
\definecolor{leanrule}{gray}{0.62}
\definecolor{leanback}{gray}{0.962}
\definecolor{leantitle}{gray}{0.22}
\definecolor{leanmuted}{gray}{0.38}
\newtcolorbox{leanbox}[1]{enhanced,breakable,colback=leanback,colframe=leanrule,
  boxrule=0.4pt,arc=0.6mm,left=2.4mm,right=2.4mm,top=1.8mm,bottom=1.8mm,
  before skip=0.8em,after skip=1em,fonttitle=\small\bfseries,coltitle=leantitle,
  colbacktitle=leanback,titlerule=0.3pt,title={#1}}
\newlength{\leancharwidth}
\newenvironment{leancode}{\par\smallskip\begingroup\ttfamily\footnotesize\raggedright
  \settowidth{\leancharwidth}{M}\setlength{\parskip}{0pt}\setlength{\parindent}{0pt}
  \frenchspacing}{\par\endgroup\smallskip}
\newcommand{\leanline}[2]{\hangindent=\dimexpr#1\leancharwidth+4\leancharwidth\relax
  \hangafter=1\noindent\hspace*{#1\leancharwidth}#2\par}
\newcommand{\leansource}[1]{\par{\footnotesize\color{leanmuted}#1}\par}
\newcommand{\leanlabel}[1]{\par\smallskip\noindent{\small\bfseries #1}\par}
\usepackage[margin=25mm]{geometry}
\usepackage[unicode,hidelinks]{hyperref}
\hypersetup{pdfauthor={LangevinLean Project}}
\setlength{\emergencystretch}{3em}
\setlength{\parskip}{0.4em}
\setlength{\parindent}{0pt}
\allowdisplaybreaks[2]
\raggedbottom
\title{\texorpdfstring{The gradient-step and descent inequalities}{The gradient-step and descent inequalities}}
\author{LangevinLean Project}
\date{}
\begin{document}\maketitle
Draft mathematical exposition; these notes may contain mistakes.


This section proves the two deterministic inequalities used in the proof of Theorem 1. They concern one gradient step \(\theta\mapsto\theta-h\nabla f(\theta)\) with step \(h\ge0\) for a differentiable potential \(f\) on a real Hilbert space \(H\) (in the applications \(H=\mathbb R^p\)). The first, the gradient-step lemma (eq:basic-gd) of the paper, uses \(m\)-strong convexity:


\[\|\theta-h\nabla f(\theta)-y\|^2\le(1-mh)\|\theta-y\|^2+2h\bigl(f(y)-f(\theta)\bigr)+h^2\|\nabla f(\theta)\|^2\qquad(\theta,y\in H).\]


The second, the deterministic half (eq:exp-upper) of the smoothing lemma, uses an \(M\)-Lipschitz gradient:


\[2\bigl(f(\theta-h\nabla f(\theta))-f(\theta)\bigr)\le-h(2-hM)\|\nabla f(\theta)\|^2.\]


In the proof of Theorem 1 the first inequality is applied with \(\theta=\vartheta\) and \(y=Y\), where \(Y\sim\pi\) is coupled with the current iterate \(\vartheta\); the second controls the potential after the gradient step. When the two are added, the potential values cancel and the coefficient \(h^2(hM-1)\) of \(\|\nabla f(\vartheta)\|^2\) remains.


\section{\texorpdfstring{The gradient-step lemma}{The gradient-step lemma}}


\textbf{Claim.} Let \(H\) be a real Hilbert space, let \(f:H\to\mathbb R\) be differentiable and \(m\)-strongly convex for some \(m\in\mathbb R\), and let \(h\ge0\). Then for all \(\theta,y\in H\),


\[\|\theta-h\nabla f(\theta)-y\|^2\le(1-mh)\|\theta-y\|^2+2h\bigl(f(y)-f(\theta)\bigr)+h^2\|\nabla f(\theta)\|^2.\]


\textbf{Proof.} Write \(g=\nabla f(\theta)\) and \(d=\theta-y\). Since \(\theta-hg-y=d-hg\), expanding the square of the norm with \(\|u-v\|^2=\|u\|^2-2\langle u,v\rangle+\|v\|^2\) and \(\|hg\|=|h|\,\|g\|=h\|g\|\) gives


\[\|\theta-hg-y\|^2=\|d\|^2-2h\langle d,g\rangle+h^2\|g\|^2.\qquad(1)\]


The supporting-plane inequality of a differentiable \(m\)-strongly convex function (inequality (3) of the previous section, valid in every real Hilbert space) at the base point \(\theta\) and the point \(y\) reads


\[f(\theta)+\langle g,y-\theta\rangle+\frac m2\|y-\theta\|^2\le f(y).\]


Since \(y-\theta=-d\), we have \(\langle g,y-\theta\rangle=-\langle d,g\rangle\) and \(\|y-\theta\|=\|d\|\), so the inequality becomes


\[-\langle d,g\rangle\le f(y)-f(\theta)-\frac m2\|d\|^2.\qquad(2)\]


Multiplying (2) by the nonnegative number \(2h\) preserves the inequality: \(-2h\langle d,g\rangle\le2h(f(y)-f(\theta))-mh\|d\|^2\). Substituting this bound into (1),


\[\|\theta-hg-y\|^2\le\|d\|^2-mh\|d\|^2+2h\bigl(f(y)-f(\theta)\bigr)+h^2\|g\|^2,\]


which is the claim, since \(\|d\|^2-mh\|d\|^2=(1-mh)\|\theta-y\|^2\). \(\square\)




The hypothesis \(h\ge0\) is used twice: to write \(\|hg\|=h\|g\|\) and to multiply (2) by \(2h\). For \(h=0\) the claim reduces to \(\|\theta-y\|^2\le\|\theta-y\|^2\). No sign condition on \(m\) and no smoothness of \(\nabla f\) are needed.

\begin{leanbox}{Lean: \texttt{gradient\_\allowbreak{}step}}
\leansource{\texttt{gradient\_\allowbreak{}step} \quad AverageLMC/\allowbreak{}GradientStep.lean, lines 24–39}
\begin{leancode}
\leanline{0}{theorem\ gradient\_step\ (f\ :\ H\ →\ ℝ)\ (m\ h\ :\ ℝ)\ (hdiff\ :\ Differentiable\ ℝ\ f)}
\leanline{4}{(hsc\ :\ StrongConvexOn\ Set.univ\ m\ f)\ (hh\ :\ 0\ ≤\ h)\ (θ\ y\ :\ H)\ :}
\leanline{4}{‖θ\ -\ h\ •\ gradient\ f\ θ\ -\ y‖\ \textasciicircum{}\ 2\ ≤}
\leanline{6}{(1\ -\ m\ *\ h)\ *\ ‖θ\ -\ y‖\ \textasciicircum{}\ 2\ +\ 2\ *\ h\ *\ (f\ y\ -\ f\ θ)\ +}
\leanline{8}{h\ \textasciicircum{}\ 2\ *\ ‖gradient\ f\ θ‖\ \textasciicircum{}\ 2}
\end{leancode}
\leanlabel{How the Lean proof follows the argument}
The ambient space is a type \texttt{H} with the instances \texttt{[Normed\allowbreak{}Add\allowbreak{}Comm\allowbreak{}Group H] [Inner\allowbreak{}Product\allowbreak{}Space \ensuremath{\mathbb R} H] [Complete\allowbreak{}Space H]}, that is, a real Hilbert space; \texttt{gradient f \ensuremath{\theta}} is the vector representing the derivative of \texttt{f} at \texttt{\ensuremath{\theta}} and \texttt{h \ensuremath{\bullet} v} is scalar multiplication. The hypotheses of \texttt{gradient\_\allowbreak{}step} are \texttt{hdiff : Differentiable \ensuremath{\mathbb R} f}, \texttt{hsc : Strong\allowbreak{}Convex\allowbreak{}On Set.\allowbreak{}univ m f} and \texttt{hh : 0 \ensuremath{\le} h}.



\begin{enumerate}


\item \texttt{hfirst} is the supporting-plane inequality, the library lemma \texttt{ULA.\allowbreak{}strong\_\allowbreak{}first\_\allowbreak{}order hdiff hsc \ensuremath{\theta} y}.

\item \texttt{hnorm} (\(\|y-\theta\|=\|\theta-y\|\), by \texttt{norm\_\allowbreak{}sub\_\allowbreak{}rev}) and \texttt{hinner} (\(\langle g,y-\theta\rangle=-\langle\theta-y,g\rangle\), by writing \(y-\theta=-(\theta-y)\) with \texttt{module}, then \texttt{inner\_\allowbreak{}neg\_\allowbreak{}right} and \texttt{real\_\allowbreak{}inner\_\allowbreak{}comm}) rewrite \texttt{hfirst} into (2).

\item \texttt{hscaled} multiplies (2) by \(2h\ge0\) with \texttt{mul\_\allowbreak{}le\_\allowbreak{}mul\_\allowbreak{}of\_\allowbreak{}nonneg\_\allowbreak{}left}; the nonnegativity of \(2h\) is checked by \texttt{positivity}.

\item The goal is rewritten into the expansion (1): the vector identity \(\theta-hg-y=(\theta-y)-hg\) (by \texttt{module}), \texttt{norm\_\allowbreak{}sub\_\allowbreak{}sq\_\allowbreak{}real} for the square of a difference, \texttt{real\_\allowbreak{}inner\_\allowbreak{}smul\_\allowbreak{}right} to extract \(h\) from the inner product, and \texttt{norm\_\allowbreak{}smul}, \texttt{Real.\allowbreak{}norm\_\allowbreak{}eq\_\allowbreak{}abs}, \texttt{abs\_\allowbreak{}of\_\allowbreak{}nonneg hh} for \(\|hg\|=h\|g\|\). \texttt{nlinarith} then combines the rewritten goal with \texttt{hscaled}.

\end{enumerate}
\end{leanbox}


Earlier results used here: \cite[Section~3.2.2]{LangevinLean2026ULASC}.

\section{\texorpdfstring{Descent of the potential along a gradient step}{Descent of the potential along a gradient step}}


\textbf{Claim.} Let \(H\) be a real Hilbert space, let \(f:H\to\mathbb R\) be differentiable with \(M\)-Lipschitz gradient for some \(M\ge0\), and let \(\theta\in H\) and \(h\ge0\). Then


\[2\bigl(f(\theta-h\nabla f(\theta))-f(\theta)\bigr)\le-h(2-hM)\|\nabla f(\theta)\|^2.\qquad(3)\]


If moreover \(hM\le2\), then \(f(\theta-h\nabla f(\theta))\le f(\theta)\).




\textbf{Proof.} Write \(g=\nabla f(\theta)\). The descent inequality \(f(z)\le f(\theta)+\langle g,z-\theta\rangle+\frac M2\|z-\theta\|^2\), proved in the previous section from the fundamental theorem of calculus along the segment from \(\theta\) to \(z\), applies to \(z=\theta-hg\). There \(z-\theta=-hg\), so \(\langle g,z-\theta\rangle=-h\langle g,g\rangle=-h\|g\|^2\) and \(\|z-\theta\|^2=|-h|^2\|g\|^2=h^2\|g\|^2\). Hence


\[f(\theta-hg)\le f(\theta)-h\|g\|^2+\frac M2h^2\|g\|^2.\]


Subtracting \(f(\theta)\) and multiplying by \(2\),


\[2\bigl(f(\theta-hg)-f(\theta)\bigr)\le-2h\|g\|^2+Mh^2\|g\|^2=-h(2-hM)\|g\|^2,\]


which is (3). If \(hM\le2\), the three factors \(-h\le0\), \(2-hM\ge0\) and \(\|g\|^2\ge0\) make the right-hand side of (3) a product of a nonpositive and two nonnegative numbers, hence \(\le0\). Therefore \(2(f(\theta-hg)-f(\theta))\le0\), which is the second assertion. \(\square\)




In Theorem 1 the step satisfies \(hM_\infty\le1\), so the second assertion applies: the gradient step does not increase the potential. Only (3) itself enters the one-step bound, where its right-hand side combines with the term \(h^2\|\nabla f(\vartheta)\|^2\) of the gradient-step lemma.

\begin{leanbox}{Lean: \texttt{gradient\_\allowbreak{}descent}, \texttt{gradient\_\allowbreak{}descent\_\allowbreak{}nonincreasing}}
\leansource{\texttt{gradient\_\allowbreak{}descent} \quad AverageLMC/\allowbreak{}GradientStep.lean, lines 43–51}
\begin{leancode}
\leanline{0}{theorem\ gradient\_descent\ (f\ :\ H\ →\ ℝ)\ (M\ :\ ℝ≥0)\ (hdiff\ :\ Differentiable\ ℝ\ f)}
\leanline{4}{(hLip\ :\ LipschitzWith\ M\ (gradient\ f))\ (θ\ :\ H)\ (h\ :\ ℝ)\ (hh\ :\ 0\ ≤\ h)\ :}
\leanline{4}{2\ *\ (f\ (θ\ -\ h\ •\ gradient\ f\ θ)\ -\ f\ θ)\ ≤}
\leanline{6}{-h\ *\ (2\ -\ h\ *\ (M\ :\ ℝ))\ *\ ‖gradient\ f\ θ‖\ \textasciicircum{}\ 2}
\end{leancode}
\leansource{\texttt{gradient\_\allowbreak{}descent\_\allowbreak{}nonincreasing} \quad AverageLMC/\allowbreak{}GradientStep.lean, lines 54–63}
\begin{leancode}
\leanline{0}{theorem\ gradient\_descent\_nonincreasing\ (f\ :\ H\ →\ ℝ)\ (M\ :\ ℝ≥0)}
\leanline{4}{(hdiff\ :\ Differentiable\ ℝ\ f)\ (hLip\ :\ LipschitzWith\ M\ (gradient\ f))}
\leanline{4}{(θ\ :\ H)\ (h\ :\ ℝ)\ (hh\ :\ 0\ ≤\ h)\ (hstep\ :\ h\ *\ (M\ :\ ℝ)\ ≤\ 2)\ :}
\leanline{4}{f\ (θ\ -\ h\ •\ gradient\ f\ θ)\ ≤\ f\ θ}
\end{leancode}
\leanlabel{How the Lean proof follows the argument}
The Lipschitz hypothesis is \texttt{h\allowbreak{}Lip : Lipschitz\allowbreak{}With M (gradient f)} with \texttt{M : \ensuremath{\mathbb R}\ensuremath{\ge}0}, a nonnegative real; \texttt{(M : \ensuremath{\mathbb R})} is its coercion to \(\mathbb R\).



\begin{enumerate}


\item In \texttt{gradient\_\allowbreak{}descent}, \texttt{hs} is the descent inequality \texttt{ULA.\allowbreak{}Analysis.\allowbreak{}smooth\_\allowbreak{}upper\_\allowbreak{}bound} applied at \texttt{\ensuremath{\theta}} and \texttt{\ensuremath{\theta} -\allowbreak{} h \ensuremath{\bullet} gradient f \ensuremath{\theta}}. \texttt{heq} computes the displacement \(-hg\) (by \texttt{module}), and the rewrites \texttt{real\_\allowbreak{}inner\_\allowbreak{}smul\_\allowbreak{}right}, \texttt{real\_\allowbreak{}inner\_\allowbreak{}self\_\allowbreak{}eq\_\allowbreak{}norm\_\allowbreak{}sq}, \texttt{norm\_\allowbreak{}smul}, \texttt{Real.\allowbreak{}norm\_\allowbreak{}eq\_\allowbreak{}abs}, \texttt{abs\_\allowbreak{}neg} and \texttt{abs\_\allowbreak{}of\_\allowbreak{}nonneg hh} evaluate \(\langle g,-hg\rangle=-h\|g\|^2\) and \(\|-hg\|=h\|g\|\) inside \texttt{hs}. \texttt{nlinarith} multiplies by \(2\) and rearranges, giving (3).

\item \texttt{gradient\_\allowbreak{}descent\_\allowbreak{}nonincreasing} adds \texttt{hstep : h * (M : \ensuremath{\mathbb R}) \ensuremath{\le} 2}. \texttt{hneg} shows that the right-hand side of (3) is \(\le0\) by two applications of \texttt{mul\_\allowbreak{}nonpos\_\allowbreak{}of\_\allowbreak{}nonpos\_\allowbreak{}of\_\allowbreak{}nonneg}, with \texttt{neg\_\allowbreak{}nonpos.\allowbreak{}mpr hh} for \(-h\le0\), \texttt{sub\_\allowbreak{}nonneg.\allowbreak{}mpr hstep} for \(2-hM\ge0\) and \texttt{sq\_\allowbreak{}nonneg} for \(\|g\|^2\ge0\); \texttt{linarith} combines \texttt{hs} (inequality (3)) with \texttt{hneg}.

\end{enumerate}
\end{leanbox}


Earlier results used here: \cite[Section~3.9.1]{LangevinLean2026ULASC}.

\begin{thebibliography}{ULA-LSI}
\bibitem[ULA-SC]{LangevinLean2026ULASC}
LangevinLean Project. \emph{Unadjusted Langevin under Strong Convexity}. LL-01, 2026. Draft report, content revision d34ead06987c. \url{https://odakyildiz.github.io/LangevinLean/docs/pdf/UnadjustedLangevin/FullProof.pdf}.
\end{thebibliography}

\end{document}
